
\section{Exercice de synthèse : Distributeur de boissons}
On souhaite écrire un programme simulant un petit distributeur automatique de boissons. Cet exercice réutilise les notions du TP1 (variables, saisies \texttt{get\_string}/\texttt{get\_int}/\texttt{get\_float}, opérateurs arithmétiques, gestion d'erreurs) ainsi que la structure \texttt{switch case} vue au TP précédent.

\textbf{Durée estimée :} 30 min

	Tout d'abord, téléchargez le squelette du TP :
	\begin{lstlisting}[language=bash,style=console]
wget https://github.com/IUT-GEII-Annecy/squelettes/releases/download/branch-2026/distributeur.zip
\end{lstlisting}

\subsection{Niveau 1 : Choix de la boisson}

\begin{UPSTIManipulation}{Distributeur de boissons}
	Le programme doit :
	\begin{enumerate}
		\item Demander le \textbf{nom} de l'utilisateur (\texttt{get\_string}) et l'accueillir : \texttt{Bonjour, <nom> !}
		\item Afficher le menu suivant :
		      \begin{center}
			      \begin{tabular}{|l|l|}
				      \hline
				      Numéro & Boisson (prix)   \\
				      \hline
				      1      & Café (0.80 €)    \\
				      2      & Thé (0.70 €)     \\
				      3      & Chocolat (1.00 €) \\
				      4      & Eau (0.50 €)     \\
				      \hline
			      \end{tabular}
		      \end{center}
		\item Demander le \textbf{numéro de boisson} choisi (\texttt{get\_int}).\\
		      \textbf{Attention :} Le choix de la boisson doit être géré par un \texttt{switch case}. Structure \texttt{if} interdite.
		\item Demander le \textbf{montant inséré} par l'utilisateur, en euros (\texttt{get\_float}).
		\item Afficher \texttt{Voici votre <boisson>.} puis \texttt{Monnaie rendue : <montant\_rendu> euros} (avec \texttt{<montant\_rendu>} = montant inséré moins le prix de la boisson, affiché avec deux décimales).\\
		      \textit{On suppose pour ce niveau que le numéro de boisson et le montant inséré sont toujours corrects.}
	\end{enumerate}

	Exemple d'exécution :
	\begin{lstlisting}[language=bash,style=console]
Entrez votre nom :
Alex
Bonjour, Alex !
1 - Cafe (0.80 euros)
2 - The (0.70 euros)
3 - Chocolat (1.00 euros)
4 - Eau (0.50 euros)
Entrez le numero de la boisson :
1
Entrez le montant insere (en euros) :
1.00
Voici votre cafe.
Monnaie rendue : 0.20 euros
	\end{lstlisting}

	\tcblower
	\begin{lstlisting}[language=bash]
check50 iut-geii-annecy/exercices/2026/info1/tp2/distributeur/niveau1
	\end{lstlisting}
\end{UPSTIManipulation}

\subsection{Niveau 2 : Gestion des erreurs}

\begin{UPSTIManipulation}{Distributeur de boissons -- Gestion des erreurs}
	Reprendre le programme du \textbf{Niveau 1}. Ajouter, avec des structures \texttt{if} (autorisées pour cette partie), la vérification suivante \textbf{avant} d'afficher le résultat de l'achat :
	\begin{itemize}
		\item Si le numéro de boisson n'est pas compris entre 1 et 4 : afficher \texttt{ERREUR : Boisson invalide} et terminer.
		\item Sinon, si le montant inséré est strictement inférieur au prix de la boisson : afficher \texttt{ERREUR : Montant insuffisant} et terminer.
		\item Sinon : comportement du Niveau 1 (achat + monnaie rendue).
	\end{itemize}

	Exemples d'exécution :
	\begin{lstlisting}[language=bash,style=console]
Entrez votre nom :
Sam
Bonjour, Sam !
1 - Cafe (0.80 euros)
2 - The (0.70 euros)
3 - Chocolat (1.00 euros)
4 - Eau (0.50 euros)
Entrez le numero de la boisson :
2
Entrez le montant insere (en euros) :
0.50
ERREUR : Montant insuffisant
	\end{lstlisting}
	\begin{lstlisting}[language=bash,style=console]
Entrez votre nom :
Sam
Bonjour, Sam !
1 - Cafe (0.80 euros)
2 - The (0.70 euros)
3 - Chocolat (1.00 euros)
4 - Eau (0.50 euros)
Entrez le numero de la boisson :
9
Entrez le montant insere (en euros) :
5.00
ERREUR : Boisson invalide
	\end{lstlisting}

	\tcblower
	\begin{lstlisting}[language=bash]
check50 iut-geii-annecy/exercices/2026/info1/tp2/distributeur/niveau2
	\end{lstlisting}
\end{UPSTIManipulation}

\section{Conclusion (A faire même si vous n'avez pas fini)}

\begin{UPSTIManipulation}{A rendre}
	\begin{itemize}
		\item[$\Box$] Placer vous dans le dossier tp2
		\item[$\Box$] Exécuter la commande suivante pour soummettre votre TP, pour en permettre le suivi.
		      \begin{lstlisting}[language=bash,style=console]
submit50 iut-geii-annecy/exercices/2026/info1/tp2/rendu
\end{lstlisting}
	\end{itemize}
\end{UPSTIManipulation}

\begin{UPSTIinfor}{Félicitations !}
	Vous venez de terminer cet exercice de synthèse. Vous savez désormais choisir entre \texttt{if} et \texttt{switch} selon la situation, et combiner ces structures avec les notions vues au TP précédent.
\end{UPSTIinfor}
